\chapter{Введение}
\label{ch:intro}



Спектральный анализ электрических сигналов является фундаментальным инструментом в современной электронике, телекоммуникациях и обработке информации,
где сложные сигналы — от аудио- и радиоволн до цифровых данных — требуют разложения на простые гармонические компоненты для оптимизации передачи, фильтрации и восстановления. 
В инженерной практике он позволяет выявлять скрытые частотные искажения, минимизировать шумы и повышать эффективность систем, таких как радары, акустические устройства или линии связи, 
но при этом сталкивается с фундаментальными ограничениями: дискретность спектров периодических сигналов, влияние длительности импульсов на ширину спектра и потеря фазовой информации, что может привести к неоднозначности интерпретации (например, при модуляции сигналов).
Для импульсных сигналов, часто используемых в цифровой технике (импульсы управления, датчики), критически важно балансировать между узкой полосой частот для экономии пропускной способности и достаточной детализацией, 
чтобы избежать размытия спектра из-за соотношений неопределённостей, где кратковременный сигнал неизбежно генерирует широкий спектр, потенциально перегружающий систему. Поэтому инженеры при проектировании устройств, таких как осциллографы или фильтры, тестируют спектры на реальных параметрах, 
чтобы оптимизировать разрешение и снизить энергозатраты без потери качества.

Одним из ключевых методов такого анализа служит преобразование Фурье — математический инструмент, разлагающий сигнал на гармоники, — реализованный в цифровых осциллографах с алгоритмом быстрого преобразования,
что позволяет в реальном времени оценивать амплитуды и частоты, учитывая эффекты дискретизации и оконных функций для минимизации артефактов.

Целью работы было изучение спектров периодических импульсных сигналов различной формы, 
анализ влияния параметров (период и длительность) на спектральные характеристики, проверка соотношений неопределённостей и ознакомление с действием RC-фильтра как простого спектрального анализатора.